\section{Discussion}
\label{sec:discussion}

\paragraph{The linear score is the limit.} Learning pays off where one weight vector must place 2
kinds of candidates or reward a rare but strong signal (added and removed seeds), but on existing
blocks the learned weights only tie equal weights, and more training data does not help. What is
missing is not data but interactions: a signal's value depends on context (the propagation rule
matters for removed seeds, content similarity for copies), which a single linear score cannot
express.

\paragraph{Learning matters more than the optimizer.} Random search with the same budget comes within 0.013 MRR
of the GA, and the GA's settings barely matter. The contribution lies more in \emph{learning} the
combination from mined labels than in the choice of optimizer; a harder search space, such as
evolving the scoring formula itself, would put the GA to a stronger test.

\paragraph{Scope and limitations.} \emph{Labels are co-changes, not necessary changes}: a commit may bundle
unrelated edits or leave a required one for later, which makes the results pessimistic rather than
optimistic. The repositories use AWS only, and thresholds and signals were designed while looking at
the same repositories that were later evaluated (the learned weights never saw their test
repositories, but our design choices did). Finally, the history is limited to the last 2,000
commits, the condition under which the model was evaluated.
